% PATCH: C-85
% Reviewer: — (final round-2 corrections; brief intake/pending/final_corrections_round2_2026-11-04.md)
% Target section: sections/results.tex lines 1-323 (Simulation Results up to Quantitative Performance
%   in Variable Topographies)
% Status: applied
%
% Applied, by audit ID (details under "### <ID>" in intake/pending/final_audit_r2_*):
% - D-17 (Diego): "does not suffer from computational deadlocks" scoped to the retained trials.
% - D-13 (Diego) + R-13: l.217 "no simpler substitute ... matches it" moderated to the
%   n = 12 result; the duplicated "same static-threshold family" clause is gone with it.
% - B.15 / D-14: l.182 scoped to Roll RMS on Rugged Terrain and Pitch RMS on the Inclined Slope (which
%   rises with severity); R-19: the restating sentence shortened.
% - B.14 / S-17: the retry no longer claims "a different dropout realization"; R-19: the closing
%   sentence of that paragraph shortened.
% - B.7 / L-35: Aversive and Obstacle paragraphs closed with a pointer to the per-level outcome.
% - B.18 / R-03: the noise caveat at the opening of sssec:noise_robustness_extra now points back to
%   the full statement instead of restating it (the figure caption keeps it, R3b-03).
% - D-47 (Diego): L_4 is the integrator's accumulator, not a ring unit.
% - D-18: "Figure ... caption" for the no-conflict ablation result (the figure does not show it).
% - R-02 = S-27: caption (b) of fig:ablation_combined keeps the numbers; the mechanistic
%   interpretation stays only in the prose; tautological clause at l.194 removed.
% - R-06: the Obstacle exclusion mentioned once in the l.130 paragraph.
% - R-20: opening of the T_response/T_switch paragraph shortened (ECDF is defined in the subcaption).
% - S-36: "full" -> "the intact circuit". S-37: "one to two update cycles" at 3.3 kHz; "threshold-based
%   basal ganglia selection logic" -> "basal-ganglia selection logic"; double space in a subcaption.
% - L-08: one panel convention in fig:variable_topo_metrics ((a)-(c); "panel A/B" for the inner panels,
%   which carry their own A./B. titles in the image).
% - S-40 #23 = L-34 = M-22: double colon in caption (c) of fig:noise_robustness_extra_combined.
% - L-41 = B-15: "zones; the line is deliberately not drawn".
% - L-21: "100.0" in cells under the "Success (%)" header. L-24: double parenthesis at l.73.
% - L-16, L-28 = M-21, L-29, L-30, L-31, B-12, B-14, S-39, L-19: Single-Stimulus, Multimodal, "stimuli of
%   different types", "on ... terrain", "the basal-ganglia mechanism", "Initially", "continuously
%   converging", "three-section", "per the stimulus-mode mapping", removed "It is important to
%   highlight", "until three successful replicates ... had been retained", U+2010 hyphen and curly
%   apostrophe replaced, "hybrid" in lower case.
%
% Not applied: M-20 (one name per locomotion mode across the paper) beyond the figure labels -- it
%   spans every section and the brief only lists the "Differential Mobile" case under A.2.
%
% Unified diff (sections/ only):
% diff --git a/sections/results.tex b/sections/results.tex
% index ebed94e..52d63a7 100644
% --- a/sections/results.tex
% +++ b/sections/results.tex
% @@ -1,7 +1,7 @@
%  \section{Results}
%  \label{sec:results}
%  
% -\revblue{Results are reported first for Gazebo simulation (single and multiple stimuli, and gait switching on the three-topology terrain map described earlier) and then for the physical platform.} 
% +\revblue{Results are reported first for Gazebo simulation (single and multiple stimuli, and gait switching on the three-section terrain map described earlier) and then for the physical platform.} 
%  
%  \subsection{Simulation Results}
%  
% @@ -33,7 +33,7 @@ Figure~\ref{fig:RESPUESTAS_EST_UNI} shows the simulated trajectories for the thr
%       \label{fig:RESPUESTAS_EST_UNI}
%  \end{figure}
%  
% -These simulated behaviors arise from the neuronal dynamics in Figure~\ref{fig:NEU_EST_UNIG}: with an obstacle on the robot's right side, the platform evades left, driven by the sensory ring, whose unit $L_{4}$ activity progressively decreases as the robot moves away.
% +These simulated behaviors arise from the neuronal dynamics in Figure~\ref{fig:NEU_EST_UNIG}: with an obstacle on the robot's right side, the platform evades left, driven by the sensory ring \revblue{and its integrator neurons, whose accumulator unit $L_{4}$} progressively decreases its activity as the robot moves away.
%  
%  \begin{figure}[t]
%      \centering
% @@ -53,7 +53,7 @@ For the appetitive scenario (target initially to the robot's right), Figure~\ref
%      \label{fig:NEU_EST_UNIB}
%  \end{figure}
%  
% -\subsubsection{Quantitative Validation: Single-Stimuli Simulation Performance}
% +\subsubsection{Quantitative Validation: Single-Stimulus Simulation Performance}
%  The performance of the platform was evaluated under isolated sensory stimuli to validate the locomotor decision network and the sensory modules individually. \revblue{In all four scenarios, three successful trials were retained at every combined-perturbation level (Figure~\ref{fig:noise_robustness_grid}); attempts discarded before reaching that count are analyzed in Section~\ref{sssec:noise_robustness_extra}.}
%  
%  \begin{figure*}[t]
% @@ -70,7 +70,7 @@ The performance of the platform was evaluated under isolated sensory stimuli to
%      \label{fig:appetitive_micro}
%  \end{figure}
%  
% -\textbf{Appetitive Targeting (see Figure~\ref{fig:noise_robustness_grid}):} The appetitive module retains task-level resilience across the full combined-perturbation sweep (Section~\ref{ssec:NoiseRobustnessEval}) (see Figure~\ref{fig:noise_robustness_grid} for the per-level $T_{sim}$/Pitch RMS/Decision Latency/$T_{switch}$ trends). \revblue{This robustness is evaluated under the combined perturbation (IMU drift, LiDAR noise, and illumination distortion injected jointly at every non-zero severity level) rather than against any single sensor's contribution in isolation, so the trends in Figure~\ref{fig:noise_robustness_grid}, across all four scenarios, cannot be attributed to any single sensing modality. Moreover, the IMU channel is perturbed only through cumulative random-walk drift: additive IMU white noise, although implemented in the simulated sensor model, was held at zero across all five severity levels (Section~\ref{ssec:NoiseRobustnessEval}).}
% +\textbf{Appetitive Targeting (see Figure~\ref{fig:noise_robustness_grid}):} The appetitive module retains task-level resilience across the full combined-perturbation sweep (Section~\ref{ssec:NoiseRobustnessEval}; per-level $T_{sim}$/Pitch RMS/Decision Latency/$T_{switch}$ trends in Figure~\ref{fig:noise_robustness_grid}). \revblue{This robustness is evaluated under the combined perturbation (IMU drift, LiDAR noise, and illumination distortion injected jointly at every non-zero severity level) rather than against any single sensor's contribution in isolation, so the trends in Figure~\ref{fig:noise_robustness_grid}, across all four scenarios, cannot be attributed to any single sensing modality. Moreover, the IMU channel is perturbed only through cumulative random-walk drift: additive IMU white noise, although implemented in the simulated sensor model, was held at zero across all five severity levels (Section~\ref{ssec:NoiseRobustnessEval}).}
%  
%  \begin{figure}[t]
%      \centering
% @@ -79,7 +79,7 @@ The performance of the platform was evaluated under isolated sensory stimuli to
%      \label{fig:aversive_micro}
%  \end{figure}
%  
% -\textbf{Aversive Escape (see Figure~\ref{fig:aversive_micro}):} The defensive repulsion properties of the architecture were re-evaluated under the combined-perturbation sweep (Figure~\ref{fig:noise_robustness_grid}). The neural firing variance exhibits a swift, step-like onset upon threat detection (structural mobilization into a defensive overdrive state), and returns to baseline once the threat leaves the field of view.
% +\textbf{Aversive Escape (see Figure~\ref{fig:aversive_micro}):} The defensive repulsion properties of the architecture were re-evaluated under the combined-perturbation sweep \revblue{(outcome per level in Figure~\ref{fig:noise_robustness_grid}; retained-trial rule in the paragraph above)}. The neural firing variance exhibits a swift, step-like onset upon threat detection (structural mobilization into a defensive overdrive state), and returns to baseline once the threat leaves the field of view.
%  
%  \begin{figure}[t]
%      \centering
% @@ -88,27 +88,27 @@ The performance of the platform was evaluated under isolated sensory stimuli to
%      \label{fig:obstacle_micro}
%  \end{figure}
%  
% -\textbf{Obstacle Evasion (see Figure~\ref{fig:obstacle_micro}):} Smooth capsule geometries for simulated obstacles prevent localized torque spikes during contact. Neural firing variance stays at a low baseline until the robot enters the obstacle's ranging threshold, at which point the avoidance sub-circuits modulate steering weights while preserving forward momentum.
% +\textbf{Obstacle Evasion (see Figure~\ref{fig:obstacle_micro}):} Smooth capsule geometries for simulated obstacles prevent localized torque spikes during contact. Neural firing variance stays at a low baseline until the robot enters the obstacle's ranging threshold, at which point the avoidance sub-circuits modulate steering weights while preserving forward momentum \revblue{(outcome per level in Figure~\ref{fig:noise_robustness_grid}; retained-trial rule in the paragraph above)}.
%  
% -\subsubsection{Response to Multiple Stimuli of Different Natures}
% +\subsubsection{Response to Multiple Stimuli of Different Types}
%  
% -For this experiment, an obstacle was placed close on the robot's right side, an aversive stimulus on the left, and an appetitive stimulus farther away, toward which the platform was expected to approach. Neuronal activities of the four control structures are shown in Figure~\ref{fig:NEU_EST_MUL}, and the resulting trajectory and behaviors in Figure~\ref{fig:RES_EST_MUL}.
% +For this experiment, an obstacle was placed close to the robot, on its right side, an aversive stimulus on the left, and an appetitive stimulus farther away, toward which the platform was expected to approach. Neuronal activities of the four control structures are shown in Figure~\ref{fig:NEU_EST_MUL}, and the resulting trajectory and behaviors in Figure~\ref{fig:RES_EST_MUL}.
%  
%  \begin{figure}[t]
%      \centering
%      \includegraphics[width=\textwidth]{NEU_EST_MUL.png}
% -    \caption{Raster plot of the neuronal activations to multiple stimuli of diverse nature in a simulation trial. \revblue{Single trial ($N=1$).}}
% +    \caption{Raster plot of the neuronal activations to multiple stimuli of different types in a simulation trial. \revblue{Single trial ($N=1$).}}
%      \label{fig:NEU_EST_MUL}
%  \end{figure}
%  
%  \begin{figure}[t]
%      \centering
%      \includegraphics[width=0.8\textwidth]{RES_EST_MUL.png}
% -    \caption{Trajectory followed in response to multiple stimuli of different natures in a simulation trial. Snapshots captured at selected time instants (0, 8, 15, 20, 25, and every 35~s thereafter until 3~min). \revblue{Single trial ($N=1$).}}
% +    \caption{Trajectory followed in response to multiple stimuli of different types in a simulation trial. Snapshots captured at selected time instants (0, 8, 15, 20, 25, and every 35~s thereafter until 3~min). \revblue{Single trial ($N=1$).}}
%      \label{fig:RES_EST_MUL}
%  \end{figure}
%  
% -At the initial moment, the evasive sensory structure (Figure~\ref{fig:NEU_EST_MUL}b) and the obstacle-associated globus pallidus activate, signaling an evasive maneuver; in the locomotion structure (Figure~\ref{fig:NEU_EST_MUL}c), units $L_0$, $L_1$, and $L_3$ activate, indicating the obstacle is on the right. As it is cleared, activity in both modules decreases.
% +Initially, the evasive sensory structure (Figure~\ref{fig:NEU_EST_MUL}b) and the obstacle-associated globus pallidus activate, signaling an evasive maneuver; in the locomotion structure (Figure~\ref{fig:NEU_EST_MUL}c), units $L_0$, $L_1$, and $L_3$ activate, indicating the obstacle is on the right. As it is cleared, activity in both modules decreases.
%  
%  At $t=21$~s, the basal ganglia module (Figure~\ref{fig:NEU_EST_MUL}a) shifts dominance to the closer aversive stimulus: the robot flees rightward, driven by units $X_4$/$X_5$, with the evasion direction set by $X_{11}$ outweighing $X_{12}$; as with the obstacle, the aversive globus pallidus activation then declines, marking the behavioral transition.
%  
% @@ -116,7 +116,7 @@ At $t=27$~s the appetitive stimulus enters view, activating its globus pallidus
%  
%  The platform reaches the goal in 92~s, sequentially adjusting locomotion mode via units $X_{14}$ (obstacle evasion, omnidirectional), $X_{15}$ (aversive escape, quadrupedal), and $X_{16}$ (appetitive approach, differential rolling). In Figure~\ref{fig:NEU_EST_MUL}d, units $X_1$/$X_2$ stay inactive (no terrain instability detected, gait unchanged), while $X_0$ shows only a small-amplitude activation producing transient states around the quadrupedal selection, attributable to that gait's abrupt movements.
%  
% -\subsubsection{Quantitative Validation: Multi-Modal Simulation Performance}
% +\subsubsection{Quantitative Validation: Multimodal Simulation Performance}
%  
%  \begin{figure}[t]
%      \centering
% @@ -125,9 +125,9 @@ The platform reaches the goal in 92~s, sequentially adjusting locomotion mode vi
%      \label{fig:complex_micro}
%  \end{figure}
%  
% -\textbf{Complex Navigation (see Figure~\ref{fig:noise_robustness_grid}):} The complex navigation suite presents the obstacle, aversive, and appetitive stimuli within a single trial. In the neural firing-variance profile of Figure~\ref{fig:complex_micro}, the robot first clears the physical obstacle through the LiDAR-driven evasion pathway, without a corresponding shaded region since this channel is not one of the two stimulus-presence traces plotted; the aversive threat then drives a red-shaded burst of firing-variance activity as the quadrupedal escape response engages; once clear of the threat, a blue-shaded interval marks sustained appetitive tracking as the platform reorients and approaches the target under differential-rolling mode. The kinematic panel's Roll/Pitch RMS response tracks this same three-stimulus sequence, rising through each gait transition and settling once the target is reached. \revblue{The arbitration network} does not suffer from computational deadlocks even when a single trial exercises all three behavioral channels under joint sensor perturbation.
% +\textbf{Complex Navigation (see Figure~\ref{fig:noise_robustness_grid}):} The complex navigation suite presents the obstacle, aversive, and appetitive stimuli within a single trial. In the neural firing-variance profile of Figure~\ref{fig:complex_micro}, the robot first clears the physical obstacle through the LiDAR-driven evasion pathway, without a corresponding shaded region since this channel is not one of the two stimulus-presence traces plotted; the aversive threat then drives a red-shaded burst of firing-variance activity as the quadrupedal escape response engages; once clear of the threat, a blue-shaded interval marks sustained appetitive tracking as the platform reorients and approaches the target under differential-rolling mode. The kinematic panel's Roll/Pitch RMS response tracks this same three-stimulus sequence, rising through each gait transition and settling once the target is reached. \revblue{In none of the retained trials did the arbitration network stall, even when a single trial exercised all three behavioral channels under joint sensor perturbation (attempts discarded before retention are analyzed in Section~\ref{sssec:noise_robustness_extra}).}
%  
% -Table~\ref{tab:consolidated_simulation_metrics} consolidates completion-time and attitude statistics across all four scenarios \revblue{for a separate archived simulation battery (15 successful trials per scenario, 14 in Obstacle Evasion after the exclusion noted below; these are the same trials as the corresponding four basal-ganglia rows of Table~\ref{tab:fsm_comparison}, whose three remaining suites come from separate batteries), not for the combined-perturbation sweep of Figure~\ref{fig:noise_robustness_grid}; its pooled values are therefore not a summary of that figure's per-level trends}. Values below are recomputed directly from the archived per-trial simulation logs (one Obstacle trial with an anomalous, physically implausible roll reading was excluded, as in Table~\ref{tab:fsm_comparison}).
% +Table~\ref{tab:consolidated_simulation_metrics} consolidates completion-time and attitude statistics across all four scenarios \revblue{for a separate archived simulation battery (15 successful trials per scenario; 14 in Obstacle Evasion, where one trial with a physically implausible roll reading was excluded, as in Table~\ref{tab:fsm_comparison}; these are the same trials as the corresponding four basal-ganglia rows of Table~\ref{tab:fsm_comparison}, whose three remaining suites come from separate batteries), not for the combined-perturbation sweep of Figure~\ref{fig:noise_robustness_grid}; its pooled values are therefore not a summary of that figure's per-level trends}. Values are recomputed directly from the archived per-trial simulation logs.
%  
%  \begin{table}[htbp]
%  \centering
% @@ -138,10 +138,10 @@ Table~\ref{tab:consolidated_simulation_metrics} consolidates completion-time and
%  \toprule
%  \textbf{Suite} & \textbf{Success (\%)} & \textbf{$T_{sim}$ (s)} & \textbf{Roll (deg)} & \textbf{Pitch (deg)} \\
%  \midrule
% -Appetitive Targeting & 100.0\% & 38.52 $\pm$ 1.20 & 0.54 $\pm$ 0.01 & 1.07 $\pm$ 0.00 \\
% -Aversive Escape      & 100.0\% & 9.98 $\pm$ 2.07  & 1.36 $\pm$ 0.11 & 1.11 $\pm$ 0.07 \\
% -Obstacle Evasion     & 100.0\% & 15.72 $\pm$ 2.49 & 1.55 $\pm$ 0.15 & 0.90 $\pm$ 0.31 \\
% -Complex Navigation   & 100.0\% & 24.96 $\pm$ 10.36 & 1.03 $\pm$ 0.35 & 1.01 $\pm$ 0.13 \\
% +Appetitive Targeting & 100.0 & 38.52 $\pm$ 1.20 & 0.54 $\pm$ 0.01 & 1.07 $\pm$ 0.00 \\
% +Aversive Escape      & 100.0 & 9.98 $\pm$ 2.07  & 1.36 $\pm$ 0.11 & 1.11 $\pm$ 0.07 \\
% +Obstacle Evasion     & 100.0 & 15.72 $\pm$ 2.49 & 1.55 $\pm$ 0.15 & 0.90 $\pm$ 0.31 \\
% +Complex Navigation   & 100.0 & 24.96 $\pm$ 10.36 & 1.03 $\pm$ 0.35 & 1.01 $\pm$ 0.13 \\
%  \bottomrule
%  \end{tabularx}
%  \end{table}
% @@ -149,7 +149,7 @@ Complex Navigation   & 100.0\% & 24.96 $\pm$ 10.36 & 1.03 $\pm$ 0.35 & 1.01 $\pm
%  \subsubsection{Combined-Perturbation Robustness: Terrain Stability, Failure Modes, and Camera Dropout}
%  \label{sssec:noise_robustness_extra}
%  
% -Two further scenarios (the Rugged Terrain course and the Inclined Slope, already introduced in Section~\ref{sssec:variable_topo_quant}) were re-evaluated under the same combined IMU-drift/LiDAR-noise/illumination-distortion sweep as the four scenarios above (Section~\ref{ssec:NoiseRobustnessEval})\revblue{, with results reported in Figure~\ref{fig:noise_robustness_extra_combined}. As in those four scenarios, the three perturbation sources are injected jointly at every non-zero severity level rather than in isolation, so the results discussed below cannot be decomposed into each sensor's individual contribution; the IMU channel, as before, is perturbed only through cumulative random-walk drift, with additive white noise held at zero throughout (Section~\ref{ssec:NoiseRobustnessEval}).} This is a separate battery from the Rugged Terrain/Inclined Slope trials analyzed in Section~\ref{sssec:variable_topo_quant}\revblue{, which carried only proportional Gaussian sensor noise,} and is not directly comparable to it; the two batteries differ in perturbation condition, not just sample size.
% +Two further scenarios (the Rugged Terrain course and the Inclined Slope, already introduced in Section~\ref{sssec:variable_topo_quant}) were re-evaluated under the same combined IMU-drift/LiDAR-noise/illumination-distortion sweep as the four scenarios above (Section~\ref{ssec:NoiseRobustnessEval})\revblue{, with results reported in Figure~\ref{fig:noise_robustness_extra_combined}. The joint-injection and drift-only caveats stated above for the four scenarios apply here unchanged.} This is a separate battery from the Rugged Terrain/Inclined Slope trials analyzed in Section~\ref{sssec:variable_topo_quant}\revblue{, which carried only proportional Gaussian sensor noise,} and is not directly comparable to it; the two batteries differ in perturbation condition, not just sample size.
%  
%  \begin{figure*}[t]
%      \centering
% @@ -173,15 +173,15 @@ Two further scenarios (the Rugged Terrain course and the Inclined Slope, already
%          \caption{Camera-dropout evidence.}
%          \label{fig:camera_dropout_evidence}
%      \end{subfigure}
% -    \caption{Combined-perturbation robustness battery\revblue{ (Section~\ref{ssec:NoiseRobustnessEval}). Experimental limitation: as in Figure~\ref{fig:noise_robustness_grid}, IMU drift, LiDAR range noise, and camera illumination distortion are injected jointly at every non-zero severity level, which precludes isolating each sensor's individual contribution to the results shown in panels (a)--(c); the IMU channel reflects cumulative random-walk drift only, with additive white noise held at zero.} (a) Trial outcome by scenario across all six suites, all severity levels pooled; each bar decomposes total attempts into successful trials retained for analysis and discarded failures by cause (timeout, tipover, crash): the Rugged Terrain suite accounts for the large majority of non-timeout failures (8 tipovers, 4 crashes out of 27 attempts), markedly more than any other suite. (b) Chassis Roll RMS and Pitch RMS versus Combined Perturbation Level (index $0$--$4$) for the Rugged Terrain and Inclined Slope suites, successful trials only; Roll RMS on Rugged Terrain remains within a $1.65^\circ$--$1.9^\circ$ band across all severity levels, Pitch RMS on the Inclined Slope ranges $2.3^\circ$--$5.7^\circ$ across the sweep (the perturbation index is left as an unconverted integer rather than a percentage, since no single percentage scale represents all three sensor modalities' differently-defined severity units). (c) Appetitive-stimulus detection fraction over time for a matched pair of Appetitive Targeting attempts under camera illumination distortion: a failed attempt (timeout) and its successful retry at the same severity level: the failed attempt's detection fraction collapses to $\approx 0.01$ for the majority of the trial, while the successful retry sustains $\approx 0.99$ detection throughout. \revblue{Simulation data.}}
% +    \caption{Combined-perturbation robustness battery\revblue{ (Section~\ref{ssec:NoiseRobustnessEval}). Experimental limitation: as in Figure~\ref{fig:noise_robustness_grid}, IMU drift, LiDAR range noise, and camera illumination distortion are injected jointly at every non-zero severity level, which precludes isolating each sensor's individual contribution to the results shown in panels (a)--(c); the IMU channel reflects cumulative random-walk drift only, with additive white noise held at zero.} (a) Trial outcome by scenario across all six suites, all severity levels pooled; each bar decomposes total attempts into successful trials retained for analysis and discarded failures by cause (timeout, tipover, crash): the Rugged Terrain suite accounts for the large majority of non-timeout failures (8 tipovers, 4 crashes out of 27 attempts), markedly more than any other suite. (b) Chassis Roll RMS and Pitch RMS versus Combined Perturbation Level (index $0$--$4$) for the Rugged Terrain and Inclined Slope suites, successful trials only; Roll RMS on Rugged Terrain remains within a $1.65^\circ$--$1.9^\circ$ band across all severity levels, Pitch RMS on the Inclined Slope ranges $2.3^\circ$--$5.7^\circ$ across the sweep (the perturbation index is left as an unconverted integer rather than a percentage, since no single percentage scale represents all three sensor modalities' differently-defined severity units). (c) Appetitive-stimulus detection fraction over time for a matched pair of Appetitive Targeting attempts under camera illumination distortion (a failed attempt, timeout, and its successful retry at the same severity level). The failed attempt's detection fraction collapses to $\approx 0.01$ for the majority of the trial, while the successful retry sustains $\approx 0.99$ detection throughout. \revblue{Simulation data.}}
%      \label{fig:noise_robustness_extra_combined}
%  \end{figure*}
%  
% -Figure~\ref{fig:failure_by_family} shows that failure under combined perturbation is not evenly distributed across scenarios: the Aversive Escape and Obstacle Evasion suites completed every attempt without a single discarded trial, the Appetitive Targeting and Complex Navigation suites discarded a share of attempts to timeout (respectively, before three successful replicates per severity level were retained), and the Rugged Terrain suite alone accounts for \revblue{the tipovers and crashes (counts in the Figure~\ref{fig:noise_robustness_extra_combined} caption)}, the only suite where the robot's structural stability, not merely task completion time, was compromised under perturbation. This asymmetry is addressed explicitly in Section~\ref{ssec:Limitations}.
% +Figure~\ref{fig:failure_by_family} shows that failure under combined perturbation is not evenly distributed across scenarios: the Aversive Escape and Obstacle Evasion suites completed every attempt without a single discarded trial, the Appetitive Targeting and Complex Navigation suites discarded a share of attempts to timeout until three successful replicates per severity level had been retained, and the Rugged Terrain suite alone accounts for \revblue{the tipovers and crashes (counts in the Figure~\ref{fig:noise_robustness_extra_combined} caption)}, the only suite where the robot's structural stability, not merely task completion time, was compromised under perturbation. This asymmetry is addressed explicitly in Section~\ref{ssec:Limitations}.
%  
% -Among the trials that complete successfully, Figure~\ref{fig:terrain_stability_vs_noise} shows postural stability does not visibly degrade with perturbation severity on either terrain: Roll RMS on Rugged Terrain and Pitch RMS on the Inclined Slope both stay within the bands reported in the caption across the full severity sweep, without a monotonic trend toward instability. Read together with Figure~\ref{fig:failure_by_family}, this indicates that Rugged Terrain's elevated tipover/crash count is concentrated in the trials discarded before completion rather than reflected as progressive attitude degradation in the trials that do complete, a distinction relevant to how this suite's robustness should be characterized (Section~\ref{ssec:Limitations}).
% +\revblue{Among the trials that complete successfully, Figure~\ref{fig:terrain_stability_vs_noise} shows that Roll RMS on Rugged Terrain stays within the band reported in the caption across the full severity sweep, whereas Pitch RMS on the Inclined Slope rises with perturbation severity (range in the caption). Rugged Terrain's tipovers and crashes are therefore confined to discarded trials, not reflected as progressive roll degradation in the trials that do complete (Section~\ref{ssec:Limitations}).}
%  
% -Figure~\ref{fig:camera_dropout_evidence} isolates a case where a purely sensor-level event (the illumination-distortion model's probabilistic detection dropout suppressing the appetitive-stimulus detector almost entirely for one attempt) propagates all the way to a task-level failure: with the target undetected for nearly the whole trial, the appetitive channel never accumulates enough sustained activation to drive an approach, and the trial times out. The retry at the identical severity level, drawing a different dropout realization, sustains detection throughout and completes normally. This is precisely the pipeline-level (not sensor-level-only) evidence the combined-perturbation evaluation is designed to surface: the failure mode originates in the camera-illumination stage but only becomes visible as a robustness concern once its consequence for task completion is measured, not from the sensor-level distortion metric alone.
% +Figure~\ref{fig:camera_dropout_evidence} isolates a case where a purely sensor-level event (the illumination-distortion model's probabilistic detection dropout suppressing the appetitive-stimulus detector almost entirely for one attempt) propagates all the way to a task-level failure: with the target undetected for nearly the whole trial, the appetitive channel never accumulates enough sustained activation to drive an approach, and the trial times out. \revblue{The retry at the identical severity level sustains detection throughout and completes normally. This is the pipeline-level evidence the combined-perturbation evaluation is designed to surface.}
%  
%  Regarding the LiDAR-noise component of this evaluation specifically: the perturbation reaches the reactive obstacle-evasion pathway (the sensory-ring and locomotion-module units described in Section~\ref{sec:methodology}) in every scenario reported above, but, in the code that generated this dataset, does not reach the basal-ganglia arbitration competition between the obstacle, appetitive, and aversive channels, whose LiDAR-derived input channel is held at zero irrespective of sensor noise. The robustness demonstrated here for LiDAR noise should accordingly be read as established at the reactive-evasion level, not as a demonstration that the three-way arbitration competition itself tolerates LiDAR noise.
%  
% @@ -191,7 +191,7 @@ Regarding the LiDAR-noise component of this evaluation specifically: the perturb
%  
%  To isolate the contribution of the basal-ganglia arbitration network and compare it quantitatively against simple threshold-based decision logic, the four-variant ablation described in Section~\ref{sec:methodology} was evaluated across the Appetitive Targeting (no-conflict) and Complex Navigation (three-stimulus conflict) scenarios, reusing the same scenario parameters as Table~\ref{tab:consolidated_simulation_metrics} ($n=12$ trials/variant/scenario, noise level fixed at zero to isolate the structural effect of each ablation from sensory-noise robustness, which is characterized separately, earlier in this section, for the intact circuit). This study is simulation-only; no equivalent ablation was run on the physical platform.
%  
% -\revblue{Without stimulus conflict, all four variants perform statistically indistinguishably in the Appetitive Targeting scenario (Figure~\ref{fig:ablation_combined}): removing an ablation target does not break the simple case, as expected when only one channel is ever active.} Under real conflict (Complex Navigation, Figure~\ref{fig:ablation_success_rate}), a clear gradient emerges. \revblue{The two ablations that remove basal-ganglia mechanism itself, \emph{no lateral inhibition} (inhibitory connections) and \emph{no STN/STR} (neural layers), collapse the success rate significantly, which quantifies the consequence of removing inhibitory connections or key neural layers. Replacing the entire circuit with static threshold logic (\emph{threshold only}) degrades success more mildly and not significantly on this metric alone; the case against threshold logic instead rests on the continuous metrics below.}
% +\revblue{Without stimulus conflict, all four variants perform statistically indistinguishably in the Appetitive Targeting scenario (Figure~\ref{fig:ablation_combined} caption): removing an ablation target does not break the simple case, as expected when only one channel is ever active.} Under real conflict (Complex Navigation, Figure~\ref{fig:ablation_success_rate}), a clear gradient emerges. \revblue{The two ablations that remove the basal-ganglia mechanism itself, \emph{no lateral inhibition} (inhibitory connections) and \emph{no STN/STR} (neural layers), collapse the success rate significantly. Replacing the entire circuit with static threshold logic (\emph{threshold only}) degrades success more mildly and not significantly on this metric alone; the case against threshold logic instead rests on the continuous metrics below.}
%  
%  \begin{figure*}[t]
%      \centering
% @@ -208,21 +208,21 @@ To isolate the contribution of the basal-ganglia arbitration network and compare
%          \caption{Attitude stability (Roll/Pitch RMS)\revblue{: mean per variant, error bars show the 95\% confidence interval ($n=12$ trials/variant; Mann--Whitney $p$ vs.\ the intact circuit).}}
%          \label{fig:ablation_complex_stability}
%      \end{subfigure}
% -    \caption{Basal-ganglia ablation study under real stimulus conflict (Complex Navigation scenario, $n=12$ trials/variant; Fisher's exact / Mann--Whitney $p$ vs.\ the intact circuit). \revblue{The four variants are: \emph{intact circuit} (control); \emph{no lateral inhibition} (cross-channel inhibitory terms zeroed); \emph{no STN/STR} (STN and STR layers removed, stimulus routed directly into GPe/GPi); and \emph{threshold only} (circuit replaced by a static fixed-priority threshold over the raw stimulus).} (a) Success rate: the two ablations that remove basal-ganglia mechanism itself collapse success rate significantly (\emph{no lateral inhibition}: $25.0\%$, $p<0.001$; \emph{no STN/STR}: $33.3\%$, $p=0.001$); replacing the circuit with static threshold logic degrades success more mildly and not significantly at this sample size (\emph{threshold only}: $83.3\%$, $p=0.478$). Under the no-conflict Appetitive Targeting scenario (not shown), all four variants perform statistically indistinguishably ($91.7$--$100\%$, $p=1.000$ for all three comparisons). (b) Attitude stability: Roll RMS increases significantly for all three ablations relative to the intact circuit (\emph{no lateral inhibition}: $+43.6\%$, $p<0.001$; \emph{no STN/STR}: $+45.9\%$, $p<0.001$; \emph{threshold only}: $+22.0\%$, $p=0.017$). Pitch RMS shows a different pattern: it increases for \emph{threshold only} ($+6.4\%$, $p=0.0029$), consistent with the roll-instability trend, but decreases for \emph{no lateral inhibition} and \emph{no STN/STR} ($-11.3\%$ and $-5.7\%$, both $p<0.005$), plausibly reflecting reduced net locomotor activity in trials that fail outright rather than improved attitude control. \revblue{Simulation data.}}
% +    \caption{Basal-ganglia ablation study under real stimulus conflict (Complex Navigation scenario, $n=12$ trials/variant; Fisher's exact / Mann--Whitney $p$ vs.\ the intact circuit). \revblue{The four variants are: \emph{intact circuit} (control); \emph{no lateral inhibition} (cross-channel inhibitory terms zeroed); \emph{no STN/STR} (STN and STR layers removed, stimulus routed directly into GPe/GPi); and \emph{threshold only} (circuit replaced by a static fixed-priority threshold over the raw stimulus).} (a) Success rate: the two ablations that remove the basal-ganglia mechanism itself collapse success rate significantly (\emph{no lateral inhibition}: $25.0\%$, $p<0.001$; \emph{no STN/STR}: $33.3\%$, $p=0.001$); replacing the circuit with static threshold logic degrades success more mildly and not significantly at this sample size (\emph{threshold only}: $83.3\%$, $p=0.478$). Under the no-conflict Appetitive Targeting scenario (not shown), all four variants perform statistically indistinguishably ($91.7$--$100\%$, $p=1.000$ for all three comparisons). (b) Attitude stability: Roll RMS increases significantly for all three ablations relative to the intact circuit (\emph{no lateral inhibition}: $+43.6\%$, $p<0.001$; \emph{no STN/STR}: $+45.9\%$, $p<0.001$; \emph{threshold only}: $+22.0\%$, $p=0.017$). Pitch RMS shows a different pattern: it increases for \emph{threshold only} ($+6.4\%$, $p=0.0029$), consistent with the roll-instability trend, but decreases for \emph{no lateral inhibition} and \emph{no STN/STR} ($-11.3\%$ and $-5.7\%$, both $p<0.005$). \revblue{Simulation data.}}
%      \label{fig:ablation_combined}
%  \end{figure*}
%  
% -Decision latency under conflict is unaffected by removing basal-ganglia mechanism ($81.7$~ms, 95\%~CI $[56.9, 106.4]$, for \emph{no lateral inhibition}, $p=0.34$; $66.6$~ms, 95\%~CI $[41.9, 91.2]$, for \emph{no STN/STR}, $p=0.84$; vs.\ $75.9$~ms, 95\%~CI $[59.3, 92.6]$, for \emph{full}), but is significantly slower under \emph{threshold only} ($109.1$~ms, 95\%~CI $[89.2, 129.0]$, $p=0.014$), consistent with a static threshold lacking the anticipatory, continuously-converging dynamics the WTA circuit provides during conflict resolution. \revblue{Roll RMS increases significantly for all three ablations relative to the intact circuit (Figure~\ref{fig:ablation_complex_stability}), indicating that attitude destabilization under conflict is a shared consequence of weakening the arbitration mechanism, regardless of which component is removed. Pitch RMS instead diverges by ablation type: it increases for \emph{threshold only}, consistent with the roll-instability trend, but decreases for \emph{no lateral inhibition} and \emph{no STN/STR}, plausibly a symptom of reduced net locomotor activity in the majority of trials that fail outright under these two ablations rather than improved attitude control.}
% +Decision latency under conflict is unaffected by removing basal-ganglia mechanisms ($81.7$~ms, 95\%~CI $[56.9, 106.4]$, for \emph{no lateral inhibition}, $p=0.34$; $66.6$~ms, 95\%~CI $[41.9, 91.2]$, for \emph{no STN/STR}, $p=0.84$; vs.\ $75.9$~ms, 95\%~CI $[59.3, 92.6]$, for the intact circuit), but is significantly slower under \emph{threshold only} ($109.1$~ms, 95\%~CI $[89.2, 129.0]$, $p=0.014$), consistent with a static threshold lacking the anticipatory, continuously converging dynamics the WTA circuit provides during conflict resolution. \revblue{Roll RMS increases significantly for all three ablations relative to the intact circuit (Figure~\ref{fig:ablation_complex_stability}), indicating that attitude destabilization under conflict is a shared consequence of weakening the arbitration mechanism, regardless of which component is removed. Pitch RMS instead diverges by ablation type: it increases for \emph{threshold only}, consistent with the roll-instability trend, but decreases for \emph{no lateral inhibition} and \emph{no STN/STR}, plausibly a symptom of reduced net locomotor activity in the majority of trials that fail outright under these two ablations rather than improved attitude control.}
%  
% -Taken together, these results support the winner-take-all arbitration design on two separable grounds: the basal-ganglia circuit's own mechanism, not merely the presence of \emph{some} priority rule, is what resolves genuine multi-stimulus conflict; and no simpler substitute, including the same static-threshold family already used elsewhere in this architecture's Locomotion Module, matches it once continuous, real-time arbitration is what the task actually requires.
% +Taken together, these results support the winner-take-all arbitration design on two separable grounds: the basal-ganglia circuit's own mechanism, not merely the presence of \emph{some} priority rule, is what resolves genuine multi-stimulus conflict; \revblue{and the static-threshold substitute is slower and less attitude-stable, although its success rate is not statistically distinguishable at $n=12$.}
%  
%  \revblue{The study's sample size supports these two conclusions to different degrees. The first, that removing basal-ganglia mechanisms breaks conflict resolution, is robust to it; the second, that static threshold logic does not match the full circuit, is not yet conclusive. With $n=12$ simulated trials per variant and scenario, Fisher's exact test against the intact circuit's $12/12$ successes reaches $p<0.05$ only when an ablated variant fails in at least five of its twelve trials. The two mechanism-removing ablations fail in eight and nine, well beyond that boundary, so their collapse is not an artifact of the limited sample; \emph{threshold only} fails in two, so its success rate is statistically inconclusive rather than evidence of equivalence with the intact circuit, and its latency and Roll RMS effects, the closest to the significance boundary, should be regarded as provisional until replicated. Larger-scale repeated trials are left as future work to tighten these estimates and confirm the smaller effects. These would include more replicates per condition, a wider range of conflict configurations and noise levels, and an equivalent ablation protocol on the physical platform, which the present study does not include.}
%  
%  \subsubsection{Response to Variable Topographies}
%  
% -For the experiments involving variable and unstable topographies in simulation, a three‐section terrain map was designed, as shown in Figure~\ref{fig:TERR_IMU}. In the first simulated experiment, the robot begins moving forward driven by the hunting instinct encoded in the locomotor decision network. Subsequently, the transition from regular to rough terrain is introduced to analyze the standard deviation of the acceleration along the $Z$ axis throughout the simulation. \revblue{In the physical robot, this information corresponds to the MLP outputs $N_0$--$N_2$ (Section~\ref{ssec:Gait_Decision_Module}).} 
% +For the experiments involving variable and unstable topographies in simulation, a three-section terrain map was designed, as shown in Figure~\ref{fig:TERR_IMU}. In the first simulated experiment, the robot begins moving forward driven by the hunting instinct encoded in the locomotor decision network. Subsequently, the transition from regular to rough terrain is introduced to analyze the standard deviation of the acceleration along the $Z$ axis throughout the simulation. \revblue{In the physical robot, this information corresponds to the MLP outputs $N_0$--$N_2$ (Section~\ref{ssec:Gait_Decision_Module}).} 
%  
% -In a second test, the transition from flat terrain to an inclined slope is defined to evaluate the platform’s response to a sudden change in ground inclination. In this configuration, both in Gazebo and in the real robot, an aversive stimulus is positioned behind the platform; according to the previously defined behaviors, the robot must avoid it by switching to the quadrupedal mode. However, due to the slope, the platform adopts the differential rolling mode, achieving a more stable and efficient displacement. Note that while the simulated sloped terrain visually approximates a continuous curved trajectory in the Gazebo environment, the digital model (STL mesh) discretizes continuous geometry into planar facets. Consequently, the IMU telemetry in simulation captures these piecewise linear segments during motion, manifesting as alternating phases of dynamic angular change and transient stationary plateaus in the pitch response plots.
% +In a second test, the transition from flat terrain to an inclined slope is defined to evaluate the platform's response to a sudden change in ground inclination. In this configuration, both in Gazebo and in the real robot, an aversive stimulus is positioned behind the platform; per the stimulus-mode mapping, the robot must avoid it by switching to the quadrupedal mode. However, due to the slope, the platform adopts the differential rolling mode, achieving a more stable and efficient displacement. Note that while the simulated sloped terrain visually approximates a continuous curved trajectory in the Gazebo environment, the digital model (STL mesh) discretizes continuous geometry into planar facets. Consequently, the IMU telemetry in simulation captures these piecewise linear segments during motion, manifesting as alternating phases of dynamic angular change and transient stationary plateaus in the pitch response plots.
%  
%  In the first simulation experiment shown in Figure~\ref{fig:RES_IMU}, the platform transitions immediately after entering the unstable terrain, adopting an appropriate gait to maintain efficient displacement. This behavior is reflected in the neural responses (Figure~\ref{fig:NEU_IMU}), where at 16 seconds, the unit $X_{0}$ shows the highest activity due to the large variation in vertical acceleration, a direct indicator of terrain irregularities. This phenomenon is also visible in Figure~\ref{fig:GRAPH_IMU}, where the standard deviation of the robot's vertical acceleration is plotted over time. The graph highlights the instant when this value exceeds the 3.7 threshold defined in the neural dynamics.
%  
% @@ -237,7 +237,7 @@ In the first simulation experiment shown in Figure~\ref{fig:RES_IMU}, the platfo
%  
%  % GRAPH_IMU merged with GRAPH_IMU2 below, after the second-scenario prose (see fig:graph_imu_pair).
%  
% -For the second stability test in simulation, the neural activities of the gait-mode control structure were recorded (Figure~\ref{fig:NEU_IMU2}b). \revblue{From $t \approx 38$~s, neuron $X_{1}$, which responds to pitch above the $U_{Pitch}$ threshold, excites unit $X_{16}$; $X_{16}$ suppresses $X_{15}$ and commands the switch from quadrupedal to differential rolling mode.} The transition during locomotion is illustrated in Figure~\ref{fig:RES_IMU2}, with snapshots of the platform posture taken every 12 seconds from a lateral perspective. It is important to highlight the neural responses shown in Figure~\ref{fig:NEU_IMU2}a, where the persistent activation of neuron $X_{4}$ produces sustained activity in neurons $X_{11}$ and $X_{12}$, which in turn drive neuron $X_{13}$ and ultimately generate the forward escape behavior. In Figure~\ref{fig:GRAPH_IMU2}, the data recorded by the inertial sensor are shown.
% +For the second stability test in simulation, the neural activities of the gait-mode control structure were recorded (Figure~\ref{fig:NEU_IMU2}b). \revblue{From $t \approx 38$~s, neuron $X_{1}$, which responds to pitch above the $U_{Pitch}$ threshold, excites unit $X_{16}$; $X_{16}$ suppresses $X_{15}$ and commands the switch from quadrupedal to differential rolling mode.} The transition during locomotion is illustrated in Figure~\ref{fig:RES_IMU2}, with snapshots of the platform posture taken every 12 seconds from a lateral perspective. In Figure~\ref{fig:NEU_IMU2}a, the persistent activation of neuron $X_{4}$ produces sustained activity in neurons $X_{11}$ and $X_{12}$, which in turn drive neuron $X_{13}$ and ultimately generate the forward escape behavior. In Figure~\ref{fig:GRAPH_IMU2}, the data recorded by the inertial sensor are shown.
%  
%  \begin{figure}[t]
%      \centering
% @@ -251,7 +251,7 @@ For the second stability test in simulation, the neural activities of the gait-m
%    \begin{subfigure}[t]{\textwidth}
%      \centering
%      \includegraphics[width=0.85\linewidth]{RES_IMU.png}
% -    \caption{Trajectory in an unstable terrain in a simulation trial. The poses were taken every 22 seconds.}
% +    \caption{Trajectory on unstable terrain in a simulation trial. The poses were taken every 22 seconds.}
%      \label{fig:RES_IMU}
%    \end{subfigure}
%  
% @@ -272,13 +272,13 @@ For the second stability test in simulation, the neural activities of the gait-m
%    \begin{subfigure}[t]{0.48\textwidth}
%      \centering
%      \includegraphics[width=\linewidth]{GRAPH_IMU.png}
% -    \caption{Standard deviation of the acceleration in the Z axis during transport on an unstable terrain.}
% +    \caption{Standard deviation of the acceleration in the Z axis during locomotion over uneven terrain.}
%      \label{fig:GRAPH_IMU}
%    \end{subfigure}\hfill
%    \begin{subfigure}[t]{0.48\textwidth}
%      \centering
%      \includegraphics[width=\linewidth]{GRAPH_IMU2.png}
% -    \caption{Pitch angle during motion in a steeply sloping terrain.}
% +    \caption{Pitch angle during motion on steeply sloping terrain.}
%      \label{fig:GRAPH_IMU2}
%    \end{subfigure}
%    \caption{IMU signal traces motivating the mode transition in each simulated terrain scenario: (a) unstable/irregular terrain; (b) steeply sloping terrain. \revblue{Each panel shows a single trial ($N=1$ per panel).}}
% @@ -293,7 +293,7 @@ For the second stability test in simulation, the neural activities of the gait-m
%      \begin{subfigure}[t]{0.32\textwidth}
%          \centering
%          \includegraphics[width=\textwidth]{FigA_Stability_Profile.png}
% -        \caption{Dynamic stability profile (Tipover Risk, $TR$) aligned around the neural mode-switch instant ($t = 0.0\,\text{s}$) for  Rugged Terrain, upper panel and Inclined Slope, lower panel. Solid lines represent the mean across 15 independent trials; shaded bands denote $\pm 1\sigma$. The horizontal dashed line at $TR = 1.0$ marks the physical overturning boundary. No trial violated this limit across all 30 runs.}
% +        \caption{Dynamic stability profile (Tipover Risk, $TR$) aligned around the neural mode-switch instant ($t = 0.0\,\text{s}$) for Rugged Terrain (upper panel) and Inclined Slope (lower panel). Solid lines represent the mean across 15 independent trials; shaded bands denote $\pm 1\sigma$. The horizontal dashed line at $TR = 1.0$ marks the physical overturning boundary. No trial violated this limit across all 30 runs.}
%          \label{fig:stability_profile}
%      \end{subfigure}
%      \hfill
% @@ -307,19 +307,19 @@ For the second stability test in simulation, the neural activities of the gait-m
%      \begin{subfigure}[t]{0.32\textwidth}
%          \centering
%          \includegraphics[width=\textwidth]{FigC_Terrain_Adaptability.png}
% -        \caption{Kinematic phase-space scatter plot of chassis Roll RMS vs. Pitch RMS for all 30 trials. Blue squares (Rugged Terrain) and orange triangles (Inclined Slope) separate clearly along the Pitch RMS axis, with overlapping Roll RMS ranges. The dashed horizontal line marks the Pitch RMS midpoint ($\approx 1.98^{\circ}$) between the two terrain means, dividing the plot into the labeled Rugged Terrain and Inclined Slope zones, deliberately not drawn on the Roll RMS axis, where the overlap between terrains is itself the evidence of comparably effective lateral stabilization (see accompanying prose).}
% +        \caption{Kinematic phase-space scatter plot of chassis Roll RMS vs. Pitch RMS for all 30 trials. Blue squares (Rugged Terrain) and orange triangles (Inclined Slope) separate clearly along the Pitch RMS axis, with overlapping Roll RMS ranges. The dashed horizontal line marks the Pitch RMS midpoint ($\approx 1.98^{\circ}$) between the two terrain means, dividing the plot into the labeled Rugged Terrain and Inclined Slope zones; the line is deliberately not drawn on the Roll RMS axis, where the overlap between terrains is itself the evidence of comparably effective lateral stabilization (see accompanying prose).}
%          \label{fig:terrain_adaptability}
%      \end{subfigure}
% -    \caption{Quantitative performance profiles for the variable-topography battery (simulation, 15 trials per terrain): (A)~tipover risk aligned at the mode-switch instant; (B)~ECDF of mechanical transition time; (C)~Roll RMS vs.\ Pitch RMS phase space.}
% +    \caption{Quantitative performance profiles for the variable-topography battery (simulation, 15 trials per terrain): (a)~tipover risk aligned at the mode-switch instant; (b)~ECDF of mechanical transition time; (c)~Roll RMS vs.\ Pitch RMS phase space.}
%      \label{fig:variable_topo_metrics}
%  \end{figure*}
%  
% -\revblue{The Tipover Risk ($TR$) index was tracked across 15 randomized runs per terrain (Figure~\ref{fig:stability_profile}).} The $TR$ metric was formally derived from the McGhee--Frank Normalized Stability Margin as $TR = 1 - SM/r_{in}$, where $SM$ is the shortest horizontal distance between the center of mass projection and the support polygon boundary, and $r_{in}$ is the inradius of that polygon. Values approaching $TR = 0.0$ represent maximum structural safety, while $TR = 1.0$ marks the critical overturning threshold. In both terrains, the platform maintained a consistently bounded risk profile around $\mu_{TR} \approx 0.58$--$0.60$ (rugged terrain mean $0.600$; inclined slope mean $0.593$) throughout the recorded window. Rather than a monotonic escalation, the rugged-terrain trace showed an isolated transient spike reaching $TR \approx 0.77$ shortly before the mode-switch instant ($t \approx -1.3$ to $-0.5\,\text{s}$), consistent with the periodic contraction of the support polygon as individual legs were lifted during the quadrupedal gait cycle, before settling back to the $\approx 0.59$ baseline across the transition at $t = 0.0\,\text{s}$. The inclined-slope trace remained flat at $TR \approx 0.59$ across the entire window, with no escalation or drop coincident with the mode switch.
% +\revblue{The Tipover Risk ($TR$) index was tracked across 15 randomized runs per terrain (Figure~\ref{fig:stability_profile}).} The $TR$ metric was formally derived from the McGhee--Frank Normalized Stability Margin as $TR = 1 - SM/r_{in}$, where $SM$ is the shortest horizontal distance between the center of mass projection and the support polygon boundary, and $r_{in}$ is the inradius of that polygon. Values approaching $TR = 0.0$ represent maximum structural safety, while $TR = 1.0$ marks the critical overturning threshold. On both terrains, the platform maintained a consistently bounded risk profile around $\mu_{TR} \approx 0.58$--$0.60$ (rugged terrain mean $0.600$; inclined slope mean $0.593$) throughout the recorded window. Rather than a monotonic escalation, the rugged-terrain trace showed an isolated transient spike reaching $TR \approx 0.77$ shortly before the mode-switch instant ($t \approx -1.3$ to $-0.5\,\text{s}$), consistent with the periodic contraction of the support polygon as individual legs were lifted during the quadrupedal gait cycle, before settling back to the $\approx 0.59$ baseline across the transition at $t = 0.0\,\text{s}$. The inclined-slope trace remained flat at $TR \approx 0.59$ across the entire window, with no escalation or drop coincident with the mode switch.
%  
%  
% -To quantify the computational efficiency and temporal determinism of the proposed bio-inspired control architecture under stochastic conditions, the neural response latency ($T_{\text{response}}$) and the Empirical Cumulative Distribution Function (ECDF) of the mechanical state-switching delay ($T_{\text{switch}}$) were analyzed across 15 randomized trials for both experimental terrains, the latter illustrated in Figure~\ref{fig:decision_delay}. \revblue{In these trials, zero-mean Gaussian noise proportional to the measured value was added to the IMU, LiDAR, and camera inputs at five levels ($\sigma = 0$, 5, 10, 20, and 30\%), with three trials per level on each terrain.} \revblue{The neural latency was} $T_{\text{response}} = 0.0\,\text{s}$ for both terrain types, showing that the threshold-based basal ganglia selection logic filtered out the injected sensor noise and dispatched the correct locomotion command within the same simulation frame as the environmental disturbance. The mechanical transition latency was likewise sub-millisecond in the vast majority of trials, but not fully invariant across terrain types. \revblue{On rugged terrain, $T_{\text{switch}}$ fell within $0.2$--$0.7\,\text{ms}$ for 13 of 15 trials (a single discrete update cycle of the ROS~2 Control loop operating at $3.3\,\text{kHz}$), while the remaining trials exhibited the long tail shown in Figure~\ref{fig:decision_delay}A, indicating that occasional trials require additional update cycles to complete the locomotion-matrix hot-swap under this terrain's contact conditions. On the inclined slope, $T_{\text{switch}}$ spread in a five-to-six-step staircase (Figure~\ref{fig:decision_delay}B) rather than collapsing onto a single value, reflecting} genuine trial-to-trial dispersion of a comparable order of magnitude rather than a perfectly deterministic delay. \revblue{Neural decision dispatch thus remained deterministic across all 30 trials, while the residual mechanical variability, including the rugged-terrain long tail, is a bounded, terrain-dependent effect rather than chattering.}
% +\revblue{$T_{\text{response}}$ and the ECDF of $T_{\text{switch}}$ (Figure~\ref{fig:decision_delay}) were analyzed over the same 15 trials per terrain.} \revblue{In these trials, zero-mean Gaussian noise proportional to the measured value was added to the IMU, LiDAR, and camera inputs at five levels ($\sigma = 0$, 5, 10, 20, and 30\%), with three trials per level on each terrain.} \revblue{The neural latency was} $T_{\text{response}} = 0.0\,\text{s}$ for both terrain types, showing that the basal-ganglia selection logic filtered out the injected sensor noise and dispatched the correct locomotion command within the same simulation frame as the environmental disturbance. The mechanical transition latency was likewise sub-millisecond in the vast majority of trials, but not fully invariant across terrain types. \revblue{On rugged terrain, $T_{\text{switch}}$ fell within $0.2$--$0.7\,\text{ms}$ for 13 of 15 trials (\revblue{one to two} update cycles of the ROS~2 Control loop operating at $3.3\,\text{kHz}$), while the remaining trials exhibited the long tail shown in Figure~\ref{fig:decision_delay}, panel~A, indicating that occasional trials require additional update cycles to complete the locomotion-matrix hot-swap under this terrain's contact conditions. On the inclined slope, $T_{\text{switch}}$ spread in a five-to-six-step staircase (Figure~\ref{fig:decision_delay}, panel~B) rather than collapsing onto a single value, reflecting} genuine trial-to-trial dispersion of a comparable order of magnitude rather than a perfectly deterministic delay. \revblue{Neural decision dispatch thus remained deterministic across all 30 trials, while the residual mechanical variability, including the rugged-terrain long tail, is a bounded, terrain-dependent effect rather than chattering.}
%  
% -\revblue{The steady-state kinematics of both terrains were mapped onto a Roll RMS ($\phi_{\text{RMS}}$) vs.\ Pitch RMS ($\theta_{\text{RMS}}$) phase space across all 30 successful trials (Figure~\ref{fig:terrain_adaptability}).} The results revealed a clear separation in Pitch RMS between the two terrains, with partial overlap in Roll RMS. The rugged terrain converged into a low-to-moderate roll, low-pitch operational zone ($\mu_{\phi} \approx 1.59^{\circ}$, $\mu_{\theta} \approx 1.67^{\circ}$), capturing the kinematic footprint of asymmetric foot-ground interactions on uneven stones while maintaining a largely level chassis. Conversely, the inclined slope clustered in a higher-pitch region with comparable roll magnitude ($\mu_{\phi} \approx 1.47^{\circ}$, $\mu_{\theta} \approx 2.29^{\circ}$). This spatial distribution reflects the geometric adaptation triggered by the Hybrid morphological transformation: during slope ascent, the chassis aligned longitudinally with the ramp gradient (higher Pitch RMS), while the deployment of the coaxial wheel base kept the lateral support polygon wide enough that roll disturbances remained similar in magnitude to the rugged-terrain baseline rather than being further suppressed. \revblue{The Pitch-RMS separation indicates that the network selects a distinct, repeatable posture for each terrain rather than applying an environment-agnostic control law.}
% +\revblue{The steady-state kinematics of both terrains were mapped onto a Roll RMS ($\phi_{\text{RMS}}$) vs.\ Pitch RMS ($\theta_{\text{RMS}}$) phase space across all 30 successful trials (Figure~\ref{fig:terrain_adaptability}).} The results revealed a clear separation in Pitch RMS between the two terrains, with partial overlap in Roll RMS. The rugged terrain converged into a low-to-moderate roll, low-pitch operational zone ($\mu_{\phi} \approx 1.59^{\circ}$, $\mu_{\theta} \approx 1.67^{\circ}$), capturing the kinematic footprint of asymmetric foot-ground interactions on uneven stones while maintaining a largely level chassis. Conversely, the inclined slope clustered in a higher-pitch region with comparable roll magnitude ($\mu_{\phi} \approx 1.47^{\circ}$, $\mu_{\theta} \approx 2.29^{\circ}$). This spatial distribution reflects the geometric adaptation triggered by the hybrid morphological transformation: during slope ascent, the chassis aligned longitudinally with the ramp gradient (higher Pitch RMS), while the deployment of the coaxial wheel base kept the lateral support polygon wide enough that roll disturbances remained similar in magnitude to the rugged-terrain baseline rather than being further suppressed. \revblue{The Pitch-RMS separation indicates that the network selects a distinct, repeatable posture for each terrain rather than applying an environment-agnostic control law.}
%  
%  \subsubsection{Quantitative Stability Analysis of Locomotion Mode Transitions}
%  \label{sssec:transition_stability}
